\documentclass{amsart}
\usepackage{amsmath,amssymb,amsthm,hyperref}
\usepackage{algorithm}
\usepackage{algpseudocode}

\newtheorem{theorem}{Theorem}
\newtheorem{lemma}{Lemma}
\newtheorem{proposition}{Proposition}
\newtheorem{corollary}{Corollary}
\newtheorem{definition}{Definition}
\newtheorem{remark}{Remark}

\DeclareMathOperator{\poly}{poly}
\newcommand{\R}{\mathbb{R}}
\newcommand{\E}{\mathbb{E}}
\renewcommand{\Pr}{\mathbb{P}}
\newcommand{\eps}{\varepsilon}

\begin{document}

\title{Data Structures for $k$-Nearest-Neighbor Search: Trade-offs and the Curse of Dimensionality}
\maketitle

\begin{abstract}
We give a self-contained treatment of the $k$-nearest-neighbor ($k$-NN) problem.
We first record the reductions that make $k$-NN equivalent (up to logarithmic
factors) to single nearest-neighbor and to the decision version of approximate
near-neighbor search.  We then prove three positive results, each with a complete
analysis: (i) for fixed dimension $d$, balanced box-decomposition trees answer
$(1+\eps)$-approximate $k$-NN queries in time $O((1/\eps)^{d}\log n + k)$ with
$O(dn)$ space; (ii) for high-dimensional $\ell_1$ data, locality-sensitive hashing
answers the $(c,r)$-approximate near-neighbor problem in time $O(n^{\rho}\,d\,\log n)$
with $O(n^{1+\rho}+nd)$ space, where $\rho\le 1/c$, which is \emph{sublinear} in $n$
with space and preprocessing polynomial in $n$ and $d$; (iii) the
Johnson--Lindenstrauss lemma reduces the ambient dimension to $O(\eps^{-2}\log n)$
while preserving all pairwise distances to within $(1\pm\eps)$.  Combining (ii) and
(iii) shows that the curse of dimensionality \emph{can} be overcome for the
approximate problem.  We finally state precisely why all known methods for the
\emph{exact} high-dimensional problem either use space exponential in $d$ or take
query time essentially linear in $n$, so that the curse is not (under SETH) removable
in the exact regime.
\end{abstract}

\section{Formalization and reductions}

Throughout, $(X,d)$ is a metric space, $P=\{p_1,\dots,p_n\}\subset X$, and a query
is a point $q\in X$.  We measure space in machine words and time in the real-RAM /
word-RAM model in which a single distance evaluation $d(q,p)$ costs $O(D)$ where
$D$ is the description length of a point ($D=d$ for $X=\R^d$, $D=O(1)$ when distances
are given by an oracle).  We write $\poly(\cdot)$ for ``bounded by a fixed
polynomial.''

\begin{definition}[$k$-NN]
A set $N_k(q)\subseteq P$ with $|N_k(q)|=k$ is a \emph{$k$-nearest-neighbor answer}
if for all $p\in N_k(q)$ and $p'\in P\setminus N_k(q)$ we have $d(q,p)\le d(q,p')$.
\end{definition}

\begin{definition}[$(1+\eps)$-approximate $k$-NN]
Let $r_k(q)$ denote the $k$-th smallest value in the multiset $\{d(q,p):p\in P\}$.
A set $A\subseteq P$ with $|A|=k$ is a \emph{$(1+\eps)$-approximate $k$-NN answer} if
$\max_{p\in A} d(q,p)\le (1+\eps)\,r_k(q)$.
\end{definition}

\begin{definition}[$(c,r)$-approximate near neighbor, $(c,r)$-ANN]
Given $c>1$ and $r>0$, the data structure must, on query $q$:
if there exists $p\in P$ with $d(q,p)\le r$, return some $p'\in P$ with
$d(q,p')\le cr$; if $d(q,p)>cr$ for all $p\in P$, return $\bot$.
(If neither premise holds, either output is acceptable.)
\end{definition}

We first reduce $k$-NN to the case $k=1$, and the approximate problem to the decision
problem $(c,r)$-ANN.  These reductions justify focusing the hard analysis on
$(c,r)$-ANN.

\begin{lemma}[$k$-NN reduces to $1$-NN]\label{lem:ktoone}
Suppose a data structure $\mathcal D$ supports, for any query $q$ and any set
$S\subseteq P$ marked ``deleted,'' the operation of returning a nearest neighbor of
$q$ among $P\setminus S$ in time $T$, with deletions costing $O(1)$ amortized.
Then exact $k$-NN is answerable in time $O(kT)$.
\end{lemma}
\begin{proof}
Run the $1$-NN query $k$ times; after each query mark the returned point deleted and
append it to the output list $A$.  After $k$ iterations $|A|=k$.  Let $a_1,\dots,a_k$
be the points in the order returned, and let $p'\in P\setminus A$.  When $a_i$ was
returned it was a nearest neighbor of $q$ in
$P\setminus\{a_1,\dots,a_{i-1}\}\supseteq\{a_i,p'\}$, hence $d(q,a_i)\le d(q,p')$.
Thus every $a_i$ satisfies $d(q,a_i)\le d(q,p')$ for all $p'\notin A$, which is the
defining property of a $k$-NN answer.  The cost is $k$ queries plus $k$ deletions,
i.e. $O(kT)$.
\end{proof}

\begin{lemma}[Approximate $k$-NN from few $(c,r)$-ANN queries]\label{lem:anntoknn}
Fix $c=1+\eps$.  Suppose for each radius $r$ in the geometric grid
$\mathcal R=\{(1+\eps)^j: j\in\mathbb Z,\ r_{\min}\le (1+\eps)^j\le r_{\max}\}$ we have
a structure that, on query $q$, reports the set $C_r(q)$ of \emph{all} points it
certifies to be within distance $cr$ of $q$, with the guarantee that $C_r(q)$ contains
every point within distance $r$ of $q$ and only points within distance $cr$.  Let
$\Delta=r_{\max}/r_{\min}$ (the spread), with $r_{\min}>0$ a lower bound on the
positive query-to-point distances.  Then a $(1+\eps)$-approximate $k$-NN answer is
computable using $O(\log_{1+\eps}\Delta)=O(\eps^{-1}\log\Delta)$ such queries.
\end{lemma}
\begin{proof}
Let $r_k=r_k(q)$.  Let $r^\star=(1+\eps)^{j^\star}$ be the smallest grid radius with
$|C_{r^\star}(q)|\ge k$.  The grid has $\lceil\log_{1+\eps}\Delta\rceil
=O(\eps^{-1}\log\Delta)$ levels; scanning levels in increasing order finds $r^\star$
within that many queries.  At the previous level $r^\star/(1+\eps)$ we had
$|C_{r^\star/(1+\eps)}(q)|<k$.  Since $C_{r^\star/(1+\eps)}(q)$ contains \emph{every}
point within distance $r^\star/(1+\eps)$, fewer than $k$ points lie within distance
$r^\star/(1+\eps)$, so $r_k>r^\star/(1+\eps)$, i.e. $r^\star<(1+\eps)r_k$.  All points
of $C_{r^\star}(q)$ lie within distance $cr^\star=(1+\eps)r^\star<(1+\eps)^2r_k$.
Select the $k$ closest to $q$ among $C_{r^\star}(q)$ by direct distance computation,
obtaining a set $A$ with $\max_{p\in A}d(q,p)\le(1+\eps)^2 r_k$.  This is valid because
$|C_{r^\star}(q)|\ge k$ and $C_{r^\star}(q)$ contains all points within $r_k\le
r^\star$ of $q$, hence contains a true $k$-NN set.  Running the whole construction with
$\eps'=\eps/3$ in place of $\eps$ gives $(1+\eps')^2\le 1+\eps$, hence a genuine
$(1+\eps)$-approximate $k$-NN answer.
\end{proof}

\begin{remark}
Lemma~\ref{lem:anntoknn} costs an $O(\eps^{-1}\log\Delta)$ factor and uses the spread.
A spread-independent reduction with $O(\log n)$ overhead also exists
(Har-Peled--Indyk--Motwani); we use the elementary spread version because it suffices
to exhibit \emph{polynomial space} and \emph{sublinear query time}, which is what the
problem asks.
\end{remark}

\section{Fixed dimension: balanced box-decomposition trees}

We give the first positive result, for $X=\R^d$ with $d$ a constant, under any
$\ell_p$ norm.  The data structure is the \emph{balanced box-decomposition} (BBD) tree
of Arya, Mount, Netanyahu, Silverman, and Wu.

\begin{definition}[BBD tree]
A \emph{box} is an axis-parallel rectangle.  A \emph{cell} is a box, possibly with a
smaller box (a ``hole'') removed.  A BBD tree on $P$ is a binary tree whose nodes are
cells, built by alternating two operations as needed: \emph{split} (cut a box by the
axis-parallel hyperplane through its midpoint orthogonal to a longest side) and
\emph{shrink} (replace a box by an inner box containing about half its points, with
the outer cell being the box minus the inner box).  Leaves contain at most one point.
\end{definition}

\begin{lemma}[BBD structural properties]\label{lem:bbd}
A BBD tree on $n$ points in $\R^d$ is constructible in $O(dn\log n)$ time, has $O(n)$
nodes, depth $O(d\log n)$, uses $O(dn)$ space, and satisfies:
\begin{enumerate}
\item[(P1)] \emph{(Balance)} along any root-to-leaf path the point count halves every
$O(d)$ levels, so the depth is $O(d\log n)$;
\item[(P2)] \emph{(Size decrease)} along any root-to-leaf path the $\ell_\infty$
diameter of the cells decreases by a factor $2$ every $O(d)$ levels;
\item[(P3)] \emph{(Packing)} for any $R\ge s>0$, the number of cells of
$\ell_\infty$-diameter $\ge s$ that intersect a ball of radius $R$ and are pairwise
disjoint is $O\big((1+R/s)^d\big)$.
\end{enumerate}
\end{lemma}
\begin{proof}
These are the standard structural properties of BBD trees; we recall the construction
giving them.  Build top-down.  At a box node, if its longest side is at most twice its
shortest side, apply \emph{split}: cut by the midpoint hyperplane orthogonal to a
longest side; after at most $d$ consecutive splits every side has been halved at least
once, giving (P2) with constant $O(d)$.  When many points concentrate in a small
sub-box (so that splits would not reduce the count quickly), apply \emph{shrink}:
compute the smallest box $b'$ containing $\lceil m/2\rceil$ of the $m$ points (found in
$O(dm)$ time by axis-wise selection), create an inner child ($b'$, $\ge m/2$ points)
and an outer child ($b\setminus b'$, $\le m/2$ points).  Alternating splits and
shrinks so that the point count halves every $O(d)$ levels gives (P1).  Hence depth is
$O(d\log n)$.  The number of nodes is $O(n)$: each shrink/split refines the partition,
and a recursion that halves the point set every $O(d)$ levels has a recursion tree with
$O(n)$ leaves and thus $O(n)$ internal nodes.  Each node stores $O(d)$ coordinates, so
space is $O(dn)$.  Construction performs $O(dn)$ work per level over $O(\log n)$
``halving'' levels, i.e. $O(dn\log n)$.  For (P3): disjoint cells of
$\ell_\infty$-diameter $\ge s$ that meet a radius-$R$ ball each contain an
$\ell_\infty$-cube of side $\ge s$ (a cell is a box minus a hole, but its outer box has
bounded aspect ratio by the split rule, so it contains such a cube up to a constant
factor); all these disjoint cubes lie in the $(R+\Theta(s))$-neighborhood of the ball,
whose $\ell_\infty$-volume is $O((R+s)^d)$, while each cube has volume $\ge(cs)^d$ for
a constant $c$; dividing gives $O((1+R/s)^d)$.
\end{proof}

\begin{theorem}[Fixed-dimension approximate $k$-NN]\label{thm:lowdim}
For $X=\R^d$ with $d$ constant, the BBD tree answers $(1+\eps)$-approximate $k$-NN
queries in time $O\big((1/\eps)^{d}\log n + k\big)$, using $O(dn)$ space and
$O(dn\log n)$ preprocessing.
\end{theorem}
\begin{proof}
We use approximate priority search (Algorithm~\ref{alg:pq}).

\begin{algorithm}[h]
\caption{Approximate priority search on a BBD tree}\label{alg:pq}
\begin{algorithmic}[1]
\Require BBD tree $\mathcal T$ on $P$, query $q$, parameter $\eps>0$, count $k$
\Ensure a $(1+\eps)$-approximate $k$-NN of $q$
\State $H\gets$ empty max-heap holding the $\le k$ best candidates, keyed by $d(q,\cdot)$
\State $Q\gets$ min-priority queue of nodes, keyed by $\mathrm{dist}(q,\mathrm{cell}(\cdot))$
\State insert the root into $Q$ with key $0$
\While{$Q\neq\emptyset$}
  \State pop node $v$ with minimum key $\delta=\mathrm{dist}(q,\mathrm{cell}(v))$
  \If{$|H|=k$ \textbf{and} $(1+\eps)\,\delta > \max H$}
     \State \textbf{break} \Comment{no unexplored cell can beat the current $k$-th best}
  \EndIf
  \If{$v$ is a leaf holding a point $p$}
     \State insert $p$ into $H$ (retain the $k$ smallest $d(q,\cdot)$)
  \Else
     \State insert both children of $v$ into $Q$ with their cell-distances to $q$
  \EndIf
\EndWhile
\State \Return the points stored in $H$
\end{algorithmic}
\end{algorithm}

Here $\mathrm{dist}(q,\mathrm{cell}(v))$ is the exact minimum $\ell_p$ distance from
$q$ to the cell, computable in $O(d)$ from the box and hole coordinates, and $\max H$
denotes the current $k$-th best (largest stored) distance.

\emph{Correctness.}  Suppose the loop breaks at line~7 with $|H|=k$ and $k$-th best
$D=\max H$.  Every node still in $Q$ has key $\ge\delta$, and every leaf reached later
holds a point $p$ with $d(q,p)\ge\delta$.  The break condition gives
$(1+\eps)\delta>D$, hence every unexplored point $p$ has
$d(q,p)\ge\delta>D/(1+\eps)$.  We claim $D\le(1+\eps)r_k$.  Let $T$ be a true $k$-NN
set.  If all of $T$ was inspected, then $D\le \max_{p\in T} d(q,p)=r_k$.  Otherwise
some $p\in T$ is unexplored, so $r_k\ge d(q,p)\ge\delta>D/(1+\eps)$, giving
$D<(1+\eps)r_k$.  In both cases $D\le(1+\eps)r_k$, the required guarantee.  If instead
$Q$ empties without breaking, all points were examined and $H$ is exact.

\emph{Running time.}  Each popped node costs $O(d)$ (cell-distance) plus $O(\log n)$
(heap/PQ operations).  We bound the number of popped nodes.  Let $R=(1+\eps)\,r_k$.
Every internal cell popped before the break has cell-distance $\le D\le R$.  Consider
the moment a popped cell $v$ has $\ell_\infty$-diameter $s_v\le \eps r_k/(2\sqrt d)$.
Then any point $p$ inside $v$ satisfies $d(q,p)\le \mathrm{dist}(q,\mathrm{cell}(v)) +
\mathrm{diam}_p(v)\le \delta+\sqrt d\,s_v\le \delta+\eps r_k/2$; if $\delta\ge r_k$,
the descendants cannot improve below $(1+\eps)r_k$, and indeed once $|H|=k$ the break
test $(1+\eps)\delta>D$ fires for $\delta>D/(1+\eps)$, terminating that branch.  Thus
the search descends into cells of diameter $\ge s:=\eps r_k/(2\sqrt d)$ that lie within
distance $R$ of $q$, and these are pairwise disjoint at any fixed scale.  By (P3) of
Lemma~\ref{lem:bbd}, the number of disjoint cells of diameter $\ge s$ within distance
$R$ is $O\big((1+R/s)^d\big)=O\big((1+ (1+\eps)r_k\cdot 2\sqrt d/(\eps r_k))^d\big)
=O\big((\sqrt d/\eps)^d\big)=O\big((1/\eps)^d\big)$ for constant $d$.  Summed over the
$O(\log n)$ scales between the root diameter and $s$ (the cell diameter halves every
$O(d)$ levels by (P2), so there are $O(\log n)$ relevant scales), the total number of
popped cells is $O\big((1/\eps)^d\log n\big)$.  Each pop costs $O(d+\log n)=O(\log n)$
for constant $d$, and reporting the heap costs $O(k)$.  Hence the query time is
$O\big((1/\eps)^d\log n+k\big)$.
\end{proof}

\begin{remark}
The exponential factor $(1/\eps)^d$ is the curse of dimensionality \emph{within this
method}: it is useless once $d=\Theta(\log n)$.  Sections~3--4 remove the exponential
$d$-dependence at the cost of approximation.
\end{remark}

\section{High dimension: locality-sensitive hashing}

Fix $c>1$, $r>0$, and consider the $(c,r)$-ANN problem.

\begin{definition}[LSH family]\label{def:lsh}
A family $\mathcal H$ of functions $h:X\to U$ is \emph{$(r,cr,p_1,p_2)$-sensitive}
(with $p_1>p_2$) if for $h$ drawn uniformly from $\mathcal H$ and any $x,y\in X$:
\[
d(x,y)\le r \ \Rightarrow\ \Pr_h[h(x)=h(y)]\ge p_1,
\qquad
d(x,y)\ge cr \ \Rightarrow\ \Pr_h[h(x)=h(y)]\le p_2 .
\]
Its \emph{quality exponent} is $\rho:=\dfrac{\ln(1/p_1)}{\ln(1/p_2)}\in(0,1)$.
\end{definition}

\begin{lemma}[Bit-sampling family for the Hamming cube]\label{lem:hamming}
Let $X=\{0,1\}^{m}$ with Hamming distance, and $\mathcal H=\{h_i(x)=x_i:1\le i\le m\}$
with $i$ uniform.  Then $\mathcal H$ is $(r,cr,p_1,p_2)$-sensitive with
$p_1=1-\frac{r}{m}$, $p_2=1-\frac{cr}{m}$, and $\rho\le 1/c$.
\end{lemma}
\begin{proof}
If $x,y$ are at Hamming distance $D$, then $\Pr_i[x_i=y_i]=1-D/m$, decreasing in $D$;
so $D\le r$ gives $\ge 1-r/m=p_1$ and $D\ge cr$ gives $\le 1-cr/m=p_2$.  For the
quality bound, $\rho\le 1/c$ is equivalent to $c\ln(1/p_1)\le\ln(1/p_2)$, i.e.
$p_2\le p_1^{\,c}$ (exponentiate $-c\ln(1/p_1)\le-\ln(1/p_2)$).  With $u=r/m\in[0,1)$,
$p_1^{\,c}=(1-u)^c$ and $p_2=1-cu$, so we need $(1-u)^c\ge 1-cu$.  Let
$g(u)=(1-u)^c-(1-cu)$; then $g(0)=0$ and $g'(u)=-c(1-u)^{c-1}+c
=c\big(1-(1-u)^{c-1}\big)\ge 0$ for $c\ge 1$ and $u\in[0,1]$.  Hence $g\ge 0$, giving
$p_2\le p_1^c$ and $\rho\le 1/c$.
\end{proof}

\begin{remark}
By the standard isometric embedding of $\ell_1$ on a bounded domain into the Hamming
cube (after scaling and discretizing each coordinate to a unary thermometer code),
distances are preserved up to a factor $1+o(1)$, so the bit-sampling family yields an
$\ell_1$ LSH family with $\rho\le 1/c$.  For $\ell_2$, the $p$-stable family
(Datar--Immorlica--Indyk) gives $\rho<1$, and Andoni--Indyk give $\rho\le
1/c^2+o(1)$.  The analysis below needs only a family with $\rho<1$ computable in $O(d)$
time; we carry $\rho$ symbolically and instantiate $\rho\le 1/c$ at the end.
\end{remark}

\begin{theorem}[LSH solves $(c,r)$-ANN in sublinear time]\label{thm:lsh}
Suppose $\mathcal H$ is $(r,cr,p_1,p_2)$-sensitive with each $h$ computable in $O(d)$
time and $p_1,p_2$ bounded away from $0$ and $1$ by constants depending only on $c$.
Let $\rho=\ln(1/p_1)/\ln(1/p_2)$ and set
\[
K:=\Big\lceil \frac{\ln n}{\ln(1/p_2)}\Big\rceil,
\qquad
L:=\Big\lceil \frac{1}{p_1}\,n^{\rho}\Big\rceil .
\]
Build $L$ hash tables; in table $\ell$, points are keyed by
$g_\ell(x)=(h_{\ell,1}(x),\dots,h_{\ell,K}(x))$ with all $h_{\ell,j}$ i.i.d.\ uniform
from $\mathcal H$.  On query $q$: scan buckets $g_1(q),\dots,g_L(q)$, compute $d(q,p)$
for each point encountered, abort after inspecting $3L$ points total, and return any
inspected $p$ with $d(q,p)\le cr$ (else $\bot$).  Then:
\begin{enumerate}
\item[(a)] Space $O(nL+nd)=O(n^{1+\rho}+nd)$; preprocessing $O(nLKd)=
O(n^{1+\rho}d\log n)$.
\item[(b)] Query time $O(LKd)=O(n^{\rho}d\log n)$.
\item[(c)] If some $p^\ast$ has $d(q,p^\ast)\le r$, the query returns a point at
distance $\le cr$ with probability $>0.29$; this is boosted to $1-\delta$ by
$O(\log(1/\delta))$ independent repetitions.  No point at distance $>cr$ is ever
returned.
\end{enumerate}
Instantiating $\rho\le 1/c$ (Lemma~\ref{lem:hamming}) gives query $O(n^{1/c}d\log n)$,
sublinear in $n$, with space and preprocessing polynomial in $n$ and $d$.
\end{theorem}
\begin{proof}
\emph{(a) Space and preprocessing.}  Each table stores the $n$ points (by reference)
under their length-$K$ keys; key strings and bucket lists cost $O(n)$ words per table
beyond the shared $O(nd)$ coordinate array, so $O(nL+nd)$ total.  Building one table
evaluates $g_\ell$ on all $n$ points, i.e. $nK$ hash evaluations at $O(d)$ each:
$O(nKd)$ per table, $O(nLKd)=O(n^{1+\rho}d\log n)$ overall, since $K=O(\log n)$ (as
$\ln(1/p_2)$ is a positive constant).

\emph{(b) Query time.}  Computing the $L$ query keys costs $O(LKd)$.  The scan
inspects at most $3L$ points, each an $O(d)$ distance, i.e. $O(Ld)$.  Total
$O(LKd)=O(n^{\rho}d\log n)$.

\emph{(c) Correctness.}  No false positive: a point is returned only after verifying
$d(q,p)\le cr$.  Now assume $p^\ast$ exists with $d(q,p^\ast)\le r$.  Define, over the
random $g_\ell$,
\[
E_1:\ \exists\,\ell\le L,\ g_\ell(p^\ast)=g_\ell(q);
\qquad
E_2:\ \sum_{\ell=1}^L \big|\{p:\,d(q,p)>cr,\ g_\ell(p)=g_\ell(q)\}\big|\le 3L .
\]
On $E_1\cap E_2$: $p^\ast$ collides with $q$ in some table, and the total number of
far collisions across all tables is $\le 3L$.  Hence among all colliding points the
far ones number $\le 3L$, so the inspection budget of $3L$ points cannot be exhausted
before $p^\ast$ (which lies in a scanned bucket) is examined; the algorithm verifies
$d(q,p^\ast)\le r\le cr$ and returns a valid point.

\emph{Bounding $\Pr[\overline{E_1}]$.}  For one hash, $\Pr[h(p^\ast)=h(q)]\ge p_1$, so
by independence $\Pr[g_\ell(p^\ast)=g_\ell(q)]\ge p_1^{\,K}$.  Let
$K_0=\ln n/\ln(1/p_2)$, so $K=\lceil K_0\rceil\in[K_0,K_0+1]$.  Since $p_1=p_2^{\rho}$
(because $\ln(1/p_1)=\rho\ln(1/p_2)$),
\[
p_1^{\,K}\ \ge\ p_1^{\,K_0+1}\ =\ p_1\cdot (p_2^{\rho})^{K_0}\ =\ p_1\cdot (p_2^{\,K_0})^{\rho}
\ =\ p_1\cdot (e^{-\ln n})^{\rho}\ =\ p_1\,n^{-\rho},
\]
using $p_2^{\,K_0}=p_2^{\,\ln n/\ln(1/p_2)}=e^{-\ln n}=1/n$.  Thus each table gives a
near collision with probability $\ge p_1 n^{-\rho}$, and with $L\ge n^{\rho}/p_1$,
\[
\Pr[\overline{E_1}]\le \big(1-p_1 n^{-\rho}\big)^{L}
\le \exp\!\big(-p_1 n^{-\rho} L\big)\le e^{-1}.
\]

\emph{Bounding $\Pr[\overline{E_2}]$.}  Fix a far point $p$, $d(q,p)>cr$.  For one
hash, $\Pr[h(p)=h(q)]\le p_2$, so $\Pr[g_\ell(p)=g_\ell(q)]\le p_2^{\,K}\le
p_2^{\,K_0}=1/n$.  Hence the expected number of far collisions in one table is
$\le n\cdot(1/n)=1$, and over $L$ tables the expected total is $\le L$.  By Markov,
$\Pr[\overline{E_2}]=\Pr[\text{total}>3L]\le 1/3$.

\emph{Combine.}  By the union bound,
$\Pr[E_1\cap E_2]\ge 1-e^{-1}-\tfrac13>0.29$.  On this event the query succeeds,
proving (c).  Running $t=\lceil \log_{1/0.71}(1/\delta)\rceil=O(\log(1/\delta))$
independent copies and returning the first success drives the failure probability to
$\le 0.71^{\,t}\le\delta$.  Instantiating $\rho\le 1/c$ gives the stated bounds and
sublinear-in-$n$ query time for every fixed $c>1$.
\end{proof}

\begin{corollary}[High-dimensional approximate $k$-NN]\label{cor:knn}
For $\ell_1$ over $\R^d$, combining Theorem~\ref{thm:lsh} (which, scanning a chosen
bucket, reports the set of certified within-$cr$ points required by
Lemma~\ref{lem:anntoknn}) with Lemma~\ref{lem:anntoknn}, $(1+\eps)$-approximate $k$-NN
is answerable with high probability in time
$O\big(n^{1/(1+\eps)}\,d\,\eps^{-1}\log n\log\Delta\big)+O(k)$ using
$O\big(n^{1+1/(1+\eps)}\eps^{-1}\log\Delta + nd\big)$ space, where $\Delta$ is the
spread.  This is sublinear in $n$ with polynomial space, for every fixed $\eps>0$.
\end{corollary}
\begin{proof}
Set $c=1+\eps$, so $\rho\le 1/(1+\eps)<1$.  For each of the $O(\eps^{-1}\log\Delta)$
grid radii build the structure of Theorem~\ref{thm:lsh}; at the level chosen by the
search of Lemma~\ref{lem:anntoknn} report all bucket points within $cr$ (their number
that we inspect is $O(n^{\rho})$ by the budget, with the same analysis).  Amplify each
level's success probability to $1-1/(n\cdot \#\text{levels})$ via $O(\log n)$ copies,
then union-bound over the $O(\eps^{-1}\log\Delta)$ levels.  Multiplying per-level cost
by the number of levels gives the bounds.
\end{proof}

\section{Dimension reduction and overcoming the curse}

\begin{theorem}[Johnson--Lindenstrauss]\label{thm:jl}
Let $\eps\in(0,1/2)$ and $V\subset\R^d$ with $|V|=N$.  Put
$m=\big\lceil 32\,\eps^{-2}\ln N\big\rceil$ and let $\Phi=\frac{1}{\sqrt m}G$ with
$G\in\R^{m\times d}$ having i.i.d.\ $\mathcal N(0,1)$ entries.  Then with probability
$\ge 1-1/N$, for all $x,y\in V$,
\[
(1-\eps)\,\|x-y\|_2^2 \le \|\Phi x-\Phi y\|_2^2 \le (1+\eps)\,\|x-y\|_2^2 .
\]
In particular such a $\Phi$ exists.
\end{theorem}
\begin{proof}
\emph{Single-vector concentration.}  Fix $v\in\R^d$, $v\ne 0$, and set
$w=v/\|v\|_2$.  The $m$ entries $Z_i:=\langle G_i,w\rangle$ (where $G_i$ is row $i$ of
$G$) are i.i.d.\ $\mathcal N(0,1)$, because each is a linear combination of independent
standard Gaussians with unit-norm coefficient vector $w$.  Then
\[
\frac{\|\Phi v\|_2^2}{\|v\|_2^2}=\frac1m\sum_{i=1}^m Z_i^2=\frac{S}{m},
\qquad S:=\sum_{i=1}^m Z_i^2\sim\chi^2_m,\quad \E[S]=m .
\]
\emph{Upper tail.}  For $0<t<\tfrac12$, $\E[e^{tZ_i^2}]=(1-2t)^{-1/2}$, so
$\E[e^{tS}]=(1-2t)^{-m/2}$.  By Markov, for any such $t$,
\[
\Pr[S\ge(1+\eps)m]\le e^{-t(1+\eps)m}(1-2t)^{-m/2}.
\]
Choosing $t=\tfrac12\cdot\tfrac{\eps}{1+\eps}\in(0,\tfrac12)$ gives $1-2t=1/(1+\eps)$
and exponent $-\tfrac m2\big[\,\eps-\ln(1+\eps)\,\big]$.  Using
$\ln(1+\eps)\le \eps-\tfrac{\eps^2}{2}+\tfrac{\eps^3}{3}$, we get
$\eps-\ln(1+\eps)\ge \tfrac{\eps^2}{2}-\tfrac{\eps^3}{3}\ge \tfrac{\eps^2}{4}$ for
$\eps\le 1/2$ (since $\tfrac{\eps^2}{2}-\tfrac{\eps^3}{3}=\eps^2(\tfrac12-\tfrac\eps3)
\ge\eps^2(\tfrac12-\tfrac16)=\tfrac{\eps^2}{3}\ge\tfrac{\eps^2}{4}$).  Hence
\[
\Pr[S\ge(1+\eps)m]\le e^{-m\eps^2/8}.
\]
\emph{Lower tail.}  For $s>0$, $\E[e^{-sZ_i^2}]=(1+2s)^{-1/2}$, so
$\E[e^{-sS}]=(1+2s)^{-m/2}$, and by Markov
\[
\Pr[S\le(1-\eps)m]=\Pr[e^{-sS}\ge e^{-s(1-\eps)m}]
\le e^{s(1-\eps)m}(1+2s)^{-m/2}.
\]
Choosing $s=\tfrac12\cdot\tfrac{\eps}{1-\eps}>0$ gives $1+2s=1/(1-\eps)$ and exponent
$-\tfrac m2\big[\,-\eps-\ln(1-\eps)\,\big]$.  Using
$-\ln(1-\eps)\ge \eps+\tfrac{\eps^2}{2}$, we get
$-\eps-\ln(1-\eps)\ge \tfrac{\eps^2}{2}\ge\tfrac{\eps^2}{4}$.  Hence
\[
\Pr[S\le(1-\eps)m]\le e^{-m\eps^2/8}.
\]
\emph{Union bound.}  For each fixed $v$,
$\Pr\big[\,|S/m-1|>\eps\,\big]\le 2e^{-m\eps^2/8}$.  Apply this to the
$\binom N2<N^2/2$ difference vectors $v=x-y$, $x\neq y\in V$.  The probability that
\emph{some} pair fails is at most
\[
\frac{N^2}{2}\cdot 2e^{-m\eps^2/8}=N^2 e^{-m\eps^2/8}.
\]
With $m=\lceil 32\eps^{-2}\ln N\rceil$ we have $m\eps^2/8\ge 4\ln N$, so
$N^2 e^{-m\eps^2/8}\le N^2 e^{-4\ln N}=N^{-2}\le 1/N$ for $N\ge 1$.  Hence with
probability $\ge 1-1/N$ every pair satisfies the two-sided bound; in particular such a
$\Phi$ exists.
\end{proof}

\begin{corollary}[Overcoming the curse for approximate NN in $\ell_2$]\label{cor:curse}
For $X=\R^d$ under $\ell_2$ and any fixed $\eps\in(0,1)$ there is a randomized data
structure for the $(1+\eps)$-approximate near-neighbor problem with:
preprocessing applies $\Phi$ of Theorem~\ref{thm:jl} (with parameter $\eps/4$,
$N=n+1$, $m=O(\eps^{-2}\log n)$) to $P$, then builds the LSH structure of
Theorem~\ref{thm:lsh} in dimension $m$; a query projects $q$ and runs the LSH query.
Its query time is
$O\big(n^{\rho}m\log n + dm\big)=O\big(n^{\rho}\eps^{-2}\log^2 n + d\,\eps^{-2}\log n\big)$
and its space is $O\big(n^{1+\rho}+nm+dm\big)$, both polynomial in $n$ and $d$, with
$\rho<1$ a constant depending only on $\eps$ (e.g.\ $\rho\le 1/(1+\eps)$ via the
$\ell_2\hookrightarrow\ell_1$ embedding, or $\rho\le 1/(1+\eps)^2+o(1)$ by
Andoni--Indyk).  Thus the query time is \emph{sublinear in $n$ while the space is
polynomial in both $n$ and $d$}: the curse of dimensionality is overcome for the
$(1+\eps)$-approximate problem.
\end{corollary}
\begin{proof}
Apply Theorem~\ref{thm:jl} to the set $V=\{q-p_i:1\le i\le n\}\cup\{0\}$ of size
$N\le n+1$, with parameter $\eps/4$.  (For online queries where $q$ is unknown at
preprocessing, use the distributional form: the single-vector concentration in the
proof of Theorem~\ref{thm:jl} holds for each fixed vector $q-p_i$ with failure
$\le 2e^{-m\eps^2/512}$, and a union bound over the $n$ vectors with
$m=O(\eps^{-2}\log n)$ yields failure $\le 1/n$ for the fixed query.)  Hence every
distance $\|q-p_i\|_2$ is preserved within a factor $(1\pm\eps/4)$ after projection.
Therefore a point that is an $(\eps/4)$-approximate near neighbor of $\Phi q$ among
$\{\Phi p_i\}$ pulls back to a point $p$ with
\[
\|q-p\|_2\le \tfrac{1}{\sqrt{1-\eps/4}}\,\|\Phi q-\Phi p\|_2
\le \tfrac{\sqrt{1+\eps/4}}{\sqrt{1-\eps/4}}\,(1+\eps/4)\,\min_i\|q-p_i\|_2
\le (1+\eps)\min_i\|q-p_i\|_2,
\]
the last inequality holding for $\eps\in(0,1)$ since
$\tfrac{(1+\eps/4)^{3/2}}{(1-\eps/4)^{1/2}}\le 1+\eps$ there.  Running
Theorem~\ref{thm:lsh} in dimension $m$ gives the stated bounds; the $d$-dependence
appears only in the projection cost $O(dm)$ per query and $O(ndm)$ in preprocessing,
which are polynomial.  The query exponent $n^{\rho}$ with constant $\rho<1$ is
sublinear in $n$.
\end{proof}

\section{The exact problem: why the curse persists}

\begin{proposition}[Exact high-dimensional barriers]\label{prop:exact}
For exact nearest-neighbor search in $\R^d$ under $\ell_2$, the two natural extreme
solutions are:
\begin{enumerate}
\item[(i)] \emph{Linear scan}: $O(nd)$ space and preprocessing, $O(nd)$ query time.
\item[(ii)] \emph{Voronoi diagram with point location}: $O(d^2\log n)$ query time but
worst-case space $n^{\Theta(d)}$.
\end{enumerate}
No known data structure simultaneously achieves query time $n^{1-\Omega(1)}$ and space
$\poly(n,d)$ for exact NN when $d=\omega(\log n)$.  Under the Strong Exponential Time
Hypothesis (SETH), no such data structure exists.
\end{proposition}
\begin{proof}[Discussion]
(i) is immediate.  For (ii), lift each $p\in P$ to
$\hat p=(p,\|p\|_2^2)\in\R^{d+1}$ and associate the hyperplane
$H_p:\,x_{d+1}=2\langle p,x_{1:d}\rangle-\|p\|_2^2$; for a query $q$,
$\|q-p\|_2^2=\|q\|_2^2-(2\langle p,q\rangle-\|p\|_2^2)$, so the nearest neighbor of $q$
corresponds to the hyperplane attaining the upper envelope at $q$, i.e. exact NN
reduces to point location in the arrangement (lower/upper envelope) of $n$ hyperplanes
in $\R^{d+1}$.  By the Upper Bound Theorem for polytopes this envelope has
$O(n^{\lceil (d+1)/2\rceil})=n^{\Theta(d)}$ faces, tight in the worst case; a
point-location structure over it answers queries in $O(d^2\log n)$ but needs space
proportional to its size, $n^{\Theta(d)}$.  The conditional lower bound is the
established fine-grained result: a data structure (equivalently, a batched algorithm
solving $n$ exact NN queries) running in $n^{2-\delta}\poly(d)$ for $d=\omega(\log n)$
would, through the reduction of CNF-SAT to Orthogonal Vectors / bichromatic closest
pair (Williams; Alman--Williams; Rubinstein), give a $2^{(1-\delta')n}$ algorithm for
CNF-SAT, contradicting SETH.  This is the precise sense in which the curse of
dimensionality holds for exact search, in contrast to the approximate case where
Corollary~\ref{cor:curse} achieves sublinear query and polynomial space.
\end{proof}

\section{Summary of the trade-offs}

\begin{itemize}
\item \textbf{Reductions.} Exact $k$-NN $\le$ $k$ calls to $1$-NN with deletions
(Lemma~\ref{lem:ktoone}); approximate $k$-NN $\le$ $O(\eps^{-1}\log\Delta)$ calls to
$(c,r)$-ANN (Lemma~\ref{lem:anntoknn}).
\item \textbf{Fixed $d$.} BBD tree: $O(dn)$ space, $O(dn\log n)$ preprocessing,
$(1+\eps)$-approximate $k$-NN in $O((1/\eps)^d\log n+k)$ (Theorem~\ref{thm:lowdim}) —
exponential in $d$, the curse in this regime.
\item \textbf{High $d$, approximate.} LSH: space $O(n^{1+\rho}+nd)$, preprocessing
$O(n^{1+\rho}d\log n)$, query $O(n^{\rho}d\log n)$, with $\rho\le 1/c$ for $\ell_1$
(and $\rho\le 1/c^2+o(1)$ for $\ell_2$): sublinear query, polynomial space
(Theorem~\ref{thm:lsh}, Corollary~\ref{cor:knn}).
\item \textbf{Dimension reduction.} JL: ambient dimension reducible to
$O(\eps^{-2}\log n)$ with $(1\pm\eps)$ distortion (Theorem~\ref{thm:jl}); composed
with LSH it overcomes the curse for $(1+\eps)$-approximate NN
(Corollary~\ref{cor:curse}).
\item \textbf{High $d$, exact.} Either $n^{\Theta(d)}$ space or $\approx n$ query time;
under SETH no $\poly(n,d)$-space, $n^{1-\Omega(1)}$-query exact algorithm exists for
$d=\omega(\log n)$ (Proposition~\ref{prop:exact}) — the curse persists.
\end{itemize}

\noindent\textbf{Conclusion.} The curse of dimensionality \emph{can} be overcome for
$(1+\eps)$-approximate nearest neighbors — sublinear query time with polynomial space,
via LSH and Johnson--Lindenstrauss — but \emph{cannot} be overcome for exact nearest
neighbors in high dimension under standard fine-grained complexity assumptions.

\end{document}
